\section{Introducción: salir a bolsa no es el punto final, sino una revisión del camino recorrido}

Esta entrevista ocurrió en un momento particular: Zhipu (智谱) confirmó que el 8 de enero de 2026 empezaría a cotizar en la Bolsa de Hong Kong, con lo que se convirtió en la primera empresa china de modelos grandes en salir a bolsa; también se le ha llamado «la primera acción de modelos grandes del mundo». Zhang Peng (张鹏) dio la entrevista con una pierna lesionada por una caída y apoyado en muletas, y el programa usa «Break a leg» como metáfora de apertura. Pero lo verdaderamente importante de este episodio no es el dramatismo del toque de campana en la bolsa, sino cómo Zhipu, a lo largo de más de seis años, pasó de la exploración de la inteligencia cognitiva en un laboratorio de Tsinghua a la industria de los modelos grandes y al gobierno corporativo de una empresa que cotiza en bolsa.

Estas notas lo organizan en tres líneas. La primera es la organizacional: cómo la tradición P2P del Laboratorio de Ingeniería del Conocimiento de Tsinghua se convirtió en la vía de transferencia de resultados de investigación de Zhipu. La segunda es la tecnológica: de la inteligencia perceptiva a la inteligencia cognitiva, sacudida una y otra vez por GPT-3, ChatGPT, DeepSeek y la Scaling Law. La tercera es la empresarial: cómo una empresa de IA de origen académico maneja el código abierto frente al código cerrado, el idealismo frente a la comercialización y el objetivo de largo plazo de la AGI frente a la disciplina de corto plazo de una empresa que cotiza en bolsa.

\begin{importantbox}{Tesis central de este episodio}
La historia de Zhipu no es la de «alcanzar de repente la ola de los modelos grandes», sino la de haber mantenido durante mucho tiempo la investigación, la ingeniería, la industria y la comercialización en una misma cadena. Salir a bolsa es solo un resultado de etapa; Zhang Peng preferiría que la historia registrara a Zhipu como «precursora de la AGI» y como «la que abrió el camino».
\end{importantbox}

\begin{knowledgebox}{Nota sobre la estrategia visual}
Este video es una toma fija de entrevista, sin slides, pizarrón ni demostraciones de producto. El cuerpo del texto no repite fotogramas de las personas; la portada sirve para identificar la fuente, y en el cuerpo el contenido didáctico se apoya en diagramas conceptuales sobre la transferencia de resultados, los sacudimientos causados por los modelos, la Scaling Law, el código abierto frente al cerrado y las etapas de la empresa.
\end{knowledgebox}

\subsection{Resumen de la sección}

Conviene leer EP129 junto con el informe trimestral sobre modelos grandes de EP136 y la historia de la tecnología de Agent de EP139. Desde la perspectiva interna de una empresa de modelos, explica cómo la industria de los modelos grandes toma forma por la acción conjunta de los paradigmas tecnológicos, los mercados de capitales, el ecosistema de código abierto y el gobierno de las organizaciones.
